% ============================================
% 第1章：数据库详细设计 + 认证权限模块详细设计
% 负责人：A
% ============================================

\section{数据库详细设计}

\subsection{数据库设计概述}

系统数据库采用PostgreSQL 16关系型数据库，数据库名称为\verb|thesis_generator|。整体遵循第三范式（3NF）设计，以论文生命周期为主线，围绕用户角色管理、模板配置管理、论文内容编辑、提交批阅流程四大业务域组织数据表结构。以下给出全部数据表的完整DDL（Data Definition Language）语句及详细字段说明。

\subsection{学院表（sys\_college）}

学院表用于存储高校下设各学院的基础信息，是整个系统数据组织体系的最上层实体。用户在注册时选择所属学院，模板按学院维度进行分类管理。

\begin{lstlisting}[language=SQL,caption={sys\_college 建表语句}]
CREATE TABLE IF NOT EXISTS sys_college (
    id BIGSERIAL PRIMARY KEY,
    name VARCHAR(100) NOT NULL,
    code VARCHAR(50) NOT NULL UNIQUE,
    created_at TIMESTAMP DEFAULT CURRENT_TIMESTAMP,
    updated_at TIMESTAMP DEFAULT CURRENT_TIMESTAMP
);
\end{lstlisting}

\begin{itemize}
    \item \verb|id|：BIGSERIAL自增主键，数据库自动分配唯一标识；
    \item \verb|name|：学院全称，VARCHAR(100)，不可为空；
    \item \verb|code|：学院编号（如CS、EE、MATH），VARCHAR(50)，唯一约束防止重复；
    \item \verb|created\_at / updated\_at|：时间戳字段，记录创建和更新时间。
\end{itemize}

\subsection{用户表（sys\_user）}

用户表是系统权限体系的核心，存储全部用户（管理员、学生、教师）的账号信息、加密密码、角色标识和学院归属。

\begin{lstlisting}[language=SQL,caption={sys\_user 建表语句}]
CREATE TABLE IF NOT EXISTS sys_user (
    id BIGSERIAL PRIMARY KEY,
    username VARCHAR(50) NOT NULL UNIQUE,
    password VARCHAR(255) NOT NULL,
    real_name VARCHAR(50) NOT NULL,
    role VARCHAR(20) NOT NULL CHECK (role IN ('ADMIN', 'STUDENT', 'TEACHER')),
    college_id BIGINT REFERENCES sys_college(id),
    enabled BOOLEAN DEFAULT TRUE,
    created_at TIMESTAMP DEFAULT CURRENT_TIMESTAMP,
    updated_at TIMESTAMP DEFAULT CURRENT_TIMESTAMP
);
\end{lstlisting}

\begin{itemize}
    \item \verb|username|：登录用户名，VARCHAR(50)，唯一约束，用于登录认证；
    \item \verb|password|：BCrypt哈希加密后的密码密文，VARCHAR(255)以容纳60字符的BCrypt哈希值；
    \item \verb|real\_name|：用户真实姓名，用于前端界面显示；
    \item \verb|role|：角色枚举，CHECK约束限定为ADMIN、STUDENT、TEACHER三种取值；
    \item \verb|college\_id|：外键引用sys\_college.id，可为空（管理员可不归属特定学院）；
    \item \verb|enabled|：账号启用状态，禁用后该账号无法登录。
\end{itemize}

角色权限体系设计如表\ref{tab:role_auth}所示：

\begin{table}[htbp]
    \centering
    \caption{角色权限矩阵}
    \label{tab:role_auth}
    \begin{tabular}{|l|c|c|c|}
        \hline
        \textbf{功能权限} & \textbf{ADMIN} & \textbf{STUDENT} & \textbf{TEACHER} \\
        \hline
        用户注册 & $\checkmark$ & $\checkmark$（仅STUDENT/TEACHER） & $\checkmark$（仅STUDENT/TEACHER） \\
        \hline
        模板管理（CRUD） & $\checkmark$ & $\times$ & $\times$ \\
        \hline
        学院管理 & $\checkmark$ & $\times$ & $\times$ \\
        \hline
        创建/编辑论文 & $\times$ & $\checkmark$ & $\times$ \\
        \hline
        文档导出 & $\times$ & $\checkmark$ & $\times$ \\
        \hline
        提交论文 & $\times$ & $\checkmark$ & $\times$ \\
        \hline
        批注/评阅论文 & $\times$ & $\times$ & $\checkmark$ \\
        \hline
        查看批阅结果 & $\times$ & $\checkmark$ & $\checkmark$ \\
        \hline
    \end{tabular}
\end{table}

\subsection{论文模板表（template）}

模板表存储由管理员创建的论文模板基本信息，每个模板可关联特定学院，并设定论文类型。

\begin{lstlisting}[language=SQL,caption={template 建表语句}]
CREATE TABLE IF NOT EXISTS template (
    id BIGSERIAL PRIMARY KEY,
    name VARCHAR(200) NOT NULL,
    type VARCHAR(20) NOT NULL CHECK (type IN ('GRADUATION', 'COURSE', 'PROJECT')),
    college_id BIGINT REFERENCES sys_college(id),
    enabled BOOLEAN DEFAULT TRUE,
    created_at TIMESTAMP DEFAULT CURRENT_TIMESTAMP,
    updated_at TIMESTAMP DEFAULT CURRENT_TIMESTAMP
);
\end{lstlisting}

\begin{itemize}
    \item \verb|name|：模板名称（如"计算机学院本科毕业论文模板"），VARCHAR(200)；
    \item \verb|type|：论文类型，CHECK约束限定三种取值：GRADUATION（毕业论文）、COURSE（课程论文）、PROJECT（项目论文）；
    \item \verb|college\_id|：外键引用sys\_college.id，实现按学院的模板隔离；
    \item \verb|enabled|：启用状态，停用的模板不显示给学生创建论文。
\end{itemize}

\subsection{模板版本表（template\_version）}

模板版本表是系统灵活性的关键设计。利用PostgreSQL特有的JSONB类型存储可变结构的配置数据，避免了为不同模板类型设计大量冗余字段的问题。

\begin{lstlisting}[language=SQL,caption={template\_version 建表语句}]
CREATE TABLE IF NOT EXISTS template_version (
    id BIGSERIAL PRIMARY KEY,
    template_id BIGINT NOT NULL REFERENCES template(id) ON DELETE CASCADE,
    version_number VARCHAR(20) NOT NULL DEFAULT '1.0',
    is_current BOOLEAN DEFAULT TRUE,
    cover_fields JSONB NOT NULL DEFAULT '[]',
    chapter_structure JSONB NOT NULL DEFAULT '[]',
    format_config JSONB NOT NULL DEFAULT '{}',
    created_at TIMESTAMP DEFAULT CURRENT_TIMESTAMP
);
\end{lstlisting}

JSONB配置字段的设计与使用场景：

\begin{itemize}
    \item \verb|cover\_fields|（JSON数组）：定义封面需填写的字段列表和显示顺序。示例：
    \begin{lstlisting}[language=Java]
[{"key":"title","label":"论文题目","required":true},
 {"key":"author","label":"作者姓名","required":true},
 {"key":"college","label":"学院","required":true}]
    \end{lstlisting}
    \item \verb|chapter\_structure|（JSON数组）：定义论文的章节树结构。示例：
    \begin{lstlisting}[language=Java]
[{"key":"abstract","title":"摘要","type":"front"},
 {"key":"ch1","title":"绪论","type":"chapter","children":[
     {"key":"ch1_1","title":"研究背景"},
     {"key":"ch1_2","title":"研究意义"}]},
 {"key":"ref","title":"参考文献","type":"back"}]
    \end{lstlisting}
    \item \verb|format\_config|（JSON对象）：定义排版格式参数。示例：
    \begin{lstlisting}[language=Java]
{"fontFamily":"SimSun","fontSize":12,"lineSpacing":1.5,
 "marginTop":2.5,"marginBottom":2.5,
 "marginLeft":3.0,"marginRight":2.5,"pageSize":"A4"}
    \end{lstlisting}
\end{itemize}

\verb|ON DELETE CASCADE|确保删除模板时其全部版本记录同步删除，维护数据一致性。

\subsection{论文主表（thesis）}

论文主表记录每篇论文的元信息和生命周期状态。

\begin{lstlisting}[language=SQL,caption={thesis 建表语句}]
CREATE TABLE IF NOT EXISTS thesis (
    id BIGSERIAL PRIMARY KEY,
    student_id BIGINT NOT NULL REFERENCES sys_user(id),
    template_version_id BIGINT REFERENCES template_version(id),
    title VARCHAR(500) NOT NULL DEFAULT '未命名论文',
    status VARCHAR(20) NOT NULL DEFAULT 'DRAFT'
        CHECK (status IN ('DRAFT','COMPLETED','SUBMITTED',
               'REVIEWED','RETURNED','GENERATING')),
    is_locked BOOLEAN DEFAULT FALSE,
    college_id BIGINT REFERENCES sys_college(id),
    created_at TIMESTAMP DEFAULT CURRENT_TIMESTAMP,
    updated_at TIMESTAMP DEFAULT CURRENT_TIMESTAMP
);
\end{lstlisting}

论文状态机流转规则如下，状态转换关系如图\ref{fig:status_flow}所示：

\begin{figure}[htbp]
    \centering
    \begin{tabular}{ccccccc}
        \fbox{DRAFT} & $\rightarrow$ & \fbox{COMPLETED} & $\rightarrow$ & \fbox{SUBMITTED} & $\rightarrow$ & \fbox{REVIEWED} \\
        & & & & \hspace{-2em}$\downarrow$ & & \\
        & & & & \fbox{RETURNED} & $\rightarrow$ & \fbox{DRAFT（解锁）} \\
    \end{tabular}
    \caption{论文状态流转图}
    \label{fig:status_flow}
\end{figure}

\begin{itemize}
    \item DRAFT → COMPLETED：学生确认论文内容完成；
    \item COMPLETED → SUBMITTED：学生提交至教师，\verb|is_locked|置为TRUE，论文只读；
    \item SUBMITTED → REVIEWED：教师完成批阅，给出分数和等级；
    \item SUBMITTED → RETURNED：教师退回论文，附带退回原因，\verb|is_locked|置为FALSE，学生可重新编辑；
    \item GENERATING：DOCX/PDF导出任务进行中的中间状态。
\end{itemize}

\subsection{论文章节表（thesis\_section）}

章节表存储每篇论文的各章节结构信息和富文本HTML内容。

\begin{lstlisting}[language=SQL,caption={thesis\_section 建表语句}]
CREATE TABLE IF NOT EXISTS thesis_section (
    id BIGSERIAL PRIMARY KEY,
    thesis_id BIGINT NOT NULL REFERENCES thesis(id) ON DELETE CASCADE,
    title VARCHAR(300) NOT NULL,
    section_key VARCHAR(100) NOT NULL,
    content TEXT DEFAULT '',
    draft_content TEXT DEFAULT '',
    section_type VARCHAR(50) DEFAULT 'chapter',
    parent_id BIGINT REFERENCES thesis_section(id),
    sort_order INT DEFAULT 0,
    created_at TIMESTAMP DEFAULT CURRENT_TIMESTAMP,
    updated_at TIMESTAMP DEFAULT CURRENT_TIMESTAMP
);
\end{lstlisting}

\begin{itemize}
    \item \verb|section\_key|：章节的业务标识键（如ch1、ch1\_1、abstract），与模板version的chapter\_structure中的key字段对应；
    \item \verb|content|：正式保存的章节HTML内容；
    \item \verb|draft\_content|：自动保存的草稿内容，独立于正式保存，30秒自动触发保存；
    \item \verb|section\_type|：章节类型标识（chapter/section/subsection/front/back），用于排版逻辑判断；
    \item \verb|parent\_id|：自引用外键，实现多级章节嵌套（父章节→子章节）；
    \item \verb|sort\_order|：同级章节的排序序号，整数递增，前端按此排序展示；
    \item \verb|ON DELETE CASCADE|：删除论文时自动删除全部章节记录。
\end{itemize}

\subsection{参考文献表（thesis\_reference）}

参考文献表遵循GB/T 7714-2015标准的数据元素设计。

\begin{lstlisting}[language=SQL,caption={thesis\_reference 建表语句}]
CREATE TABLE IF NOT EXISTS thesis_reference (
    id BIGSERIAL PRIMARY KEY,
    thesis_id BIGINT NOT NULL REFERENCES thesis(id) ON DELETE CASCADE,
    authors VARCHAR(500) NOT NULL,
    title VARCHAR(500) NOT NULL,
    journal VARCHAR(300),
    year VARCHAR(10),
    volume VARCHAR(50),
    issue VARCHAR(50),
    pages VARCHAR(50),
    doi VARCHAR(200),
    sort_order INT DEFAULT 0,
    created_at TIMESTAMP DEFAULT CURRENT_TIMESTAMP
);
\end{lstlisting}

字段设计覆盖了期刊论文、会议论文、学位论文、图书等常见文献类型的核心字段，\verb|doi|字段支持DOI链接跳转。

\subsection{图片资源表（thesis\_image）}

\begin{lstlisting}[language=SQL,caption={thesis\_image 建表语句}]
CREATE TABLE IF NOT EXISTS thesis_image (
    id BIGSERIAL PRIMARY KEY,
    thesis_id BIGINT NOT NULL REFERENCES thesis(id) ON DELETE CASCADE,
    section_id BIGINT REFERENCES thesis_section(id),
    original_name VARCHAR(500) NOT NULL,
    stored_path VARCHAR(500) NOT NULL,
    url VARCHAR(500) NOT NULL,
    file_size BIGINT DEFAULT 0,
    created_at TIMESTAMP DEFAULT CURRENT_TIMESTAMP
);
\end{lstlisting}

\begin{itemize}
    \item \verb|original\_name|：用户上传时的原始文件名；
    \item \verb|stored\_path|：服务端文件系统存储路径；
    \item \verb|url|：对外访问URL（前端img标签src引用）；
    \item \verb|file\_size|：文件大小（字节），用于空间占用统计。
\end{itemize}

\subsection{提交记录表（submission）}

\begin{lstlisting}[language=SQL,caption={submission 建表语句}]
CREATE TABLE IF NOT EXISTS submission (
    id BIGSERIAL PRIMARY KEY,
    thesis_id BIGINT NOT NULL REFERENCES thesis(id) ON DELETE CASCADE,
    student_id BIGINT NOT NULL REFERENCES sys_user(id),
    version_number INT NOT NULL DEFAULT 1,
    submitted_at TIMESTAMP DEFAULT CURRENT_TIMESTAMP
);
\end{lstlisting}

每次学生提交论文时新增一条记录，\verb|version_number|递增，保留完整的提交历史。

\subsection{批注表（annotation）}

批注表记录教师在论文指定文字位置添加的修改建议。

\begin{lstlisting}[language=SQL,caption={annotation 建表语句}]
CREATE TABLE IF NOT EXISTS annotation (
    id BIGSERIAL PRIMARY KEY,
    thesis_id BIGINT NOT NULL REFERENCES thesis(id) ON DELETE CASCADE,
    section_id BIGINT NOT NULL REFERENCES thesis_section(id),
    teacher_id BIGINT NOT NULL REFERENCES sys_user(id),
    start_offset INT NOT NULL,
    text_length INT NOT NULL,
    selected_text TEXT,
    content TEXT NOT NULL,
    created_at TIMESTAMP DEFAULT CURRENT_TIMESTAMP
);
\end{lstlisting}

\begin{itemize}
    \item \verb|start\_offset|：选中文字在原HTML内容中的起始字符位置（0-based），用于精确定位批注锚点；
    \item \verb|text\_length|：选中文字的长度，结合start\_offset确定完整范围；
    \item \verb|selected\_text|：被选中的原文内容（冗余存储，便于离线查看批注含义）；
    \item \verb|content|：教师的批注意见文本。
\end{itemize}

\subsection{批阅记录表（review\_record）}

\begin{lstlisting}[language=SQL,caption={review\_record 建表语句}]
CREATE TABLE IF NOT EXISTS review_record (
    id BIGSERIAL PRIMARY KEY,
    thesis_id BIGINT NOT NULL REFERENCES thesis(id) ON DELETE CASCADE,
    teacher_id BIGINT NOT NULL REFERENCES sys_user(id),
    comment_html TEXT,
    score INT CHECK (score >= 0 AND score <= 100),
    grade VARCHAR(10),
    action VARCHAR(20) NOT NULL CHECK (action IN ('REVIEWED', 'RETURNED')),
    return_reason TEXT,
    created_at TIMESTAMP DEFAULT CURRENT_TIMESTAMP
);
\end{lstlisting}

分数（\verb|score|）与等级（\verb|grade|）的换算规则为：$\geq$90分→优、80-89→良、70-79→中、60-69→及格、<60→不及格。\verb|action|区分批阅（REVIEWED）和退回（RETURNED）两种操作类型。

\subsection{章节草稿历史表（thesis\_section\_draft\_history）}

\begin{lstlisting}[language=SQL,caption={thesis\_section\_draft\_history 建表语句}]
CREATE TABLE IF NOT EXISTS thesis_section_draft_history (
    id BIGSERIAL PRIMARY KEY,
    section_id BIGINT NOT NULL REFERENCES thesis_section(id) ON DELETE CASCADE,
    draft_content TEXT,
    saved_at TIMESTAMP DEFAULT CURRENT_TIMESTAMP
);
\end{lstlisting}

草稿历史表用于记录每次自动保存的快照，支持学生回溯到先前的草稿版本。

\subsection{索引设计}

在概要设计中已介绍索引策略，以下给出索引的完整DDL语句：

\begin{lstlisting}[language=SQL,caption={索引创建语句}]
CREATE INDEX IF NOT EXISTS idx_user_role ON sys_user(role);
CREATE INDEX IF NOT EXISTS idx_template_college ON template(college_id);
CREATE INDEX IF NOT EXISTS idx_template_type ON template(type);
CREATE INDEX IF NOT EXISTS idx_thesis_student ON thesis(student_id);
CREATE INDEX IF NOT EXISTS idx_thesis_status ON thesis(status);
CREATE INDEX IF NOT EXISTS idx_section_thesis ON thesis_section(thesis_id);
CREATE INDEX IF NOT EXISTS idx_annotation_thesis ON annotation(thesis_id);
CREATE INDEX IF NOT EXISTS idx_review_thesis ON review_record(thesis_id);
\end{lstlisting}

全部索引均采用B-tree默认结构。由于当前数据规模为实训项目级别（预计千级用户、万级论文记录），B-tree索引足以满足所有查询场景的性能需求。若未来数据量增长至百万级，可针对JSONB字段（template\_version的配置数据）添加GIN索引以支持配置内容的高效检索。

\subsection{数据库初始化策略}

系统启动时通过Spring Boot的SQL初始化机制自动执行schema.sql脚本。在\verb|application.properties|中配置如下：

\begin{lstlisting}[language=Java,caption={SQL初始化配置}]
spring.sql.init.mode=always
spring.sql.init.schema-locations=classpath:schema.sql
spring.jpa.hibernate.ddl-auto=validate
\end{lstlisting}

\verb|ddl-auto=validate|确保JPA实体定义与实际数据库表结构保持一致，若存在字段差异则启动报错，防止运行时出现隐式的数据不一致问题。

\section{认证与权限模块详细设计}

\subsection{模块概述}

认证与权限模块（M1）负责系统中全部用户的身份认证和操作授权。模块采用JWT（JSON Web Token）无状态认证机制配合Redis缓存的自研轻量级方案，未引入Spring Security全套框架，以降低学习曲线和项目复杂度，同时满足实训项目的安全需求。

模块核心职责包括：
\begin{enumerate}
    \item \textbf{用户认证}：处理注册（创建账号）、登录（签发Token）、登出（销毁Token）三项基本认证操作；
    \item \textbf{Token管理}：JWT Token的生成、解析、过期校验，以及Redis中的缓存存储；
    \item \textbf{请求拦截}：在API请求处理管道中拦截所有非公开请求，校验Token有效性并提取用户身份信息；
    \item \textbf{角色鉴权}：根据API接口声明的所需角色，校验当前请求用户的角色是否在允许列表中。
\end{enumerate}

\subsection{认证流程设计}

\subsubsection{用户注册流程}

用户注册的完整处理流程如下：

\begin{enumerate}
    \item 前端收集用户名、密码、真实姓名、角色（STUDENT或TEACHER）、所属学院ID，通过POST请求发送至\verb|/api/v1/auth/register|；
    \item 请求体经@Valid注解触发JSR-303校验（用户名4-50位、密码6-100位、必填字段非空）；
    \item AuthController.register()接收RegisterRequest DTO，调用AuthService.register()；
    \item AuthService检查用户名是否已存在（userRepository.existsByUsername），如已存在则抛出BusinessException(400, "用户名已存在")；
    \item 校验通过后，构造User实体对象：使用BCryptPasswordEncoder.encode()对明文密码进行哈希加密，设置enabled=true，调用userRepository.save()持久化；
    \item 为新用户生成JWT Token（jwtUtil.generateToken），Token中包含userId、username、role三个Claim；
    \item 将Token写入Redis缓存（key格式：\verb|token:<userId>|），设置24小时过期，与JWT过期时间一致；
    \item 返回LoginResponse（token、userId、username、realName、role）。
\end{enumerate}

\subsubsection{用户登录流程}

登录流程与注册类似但更简化：

\begin{enumerate}
    \item 前端发送POST请求至\verb|/api/v1/auth/login|，携带LoginRequest（username + password）；
    \item AuthService根据用户名查询用户（userRepository.findByUsername），找不到则抛出BusinessException(401)；
    \item 检查enabled状态，若为false则抛出BusinessException(403, "账号已被禁用")；
    \item 使用BCryptPasswordEncoder.matches()方法比对明文密码与数据库中的密文；
    \item 认证通过后生成JWT Token并写入Redis，返回LoginResponse。
\end{enumerate}

\subsubsection{用户登出流程}

\begin{enumerate}
    \item 前端发送POST请求至\verb|/api/v1/auth/logout|（需携带Token）；
    \item AuthController从Request Attribute中提取userId（由JwtFilter预先注入）；
    \item AuthService删除Redis中对应的Token缓存；
    \item Token本身仍在其有效期内，但由于后端不再认可该Token（后续请求在Redis中找不到对应记录），实现了逻辑登出。
\end{enumerate}

\subsection{JWT Token设计}

\subsubsection{Token格式与内容}

系统采用JJWT库（io.jsonwebtoken 0.12.6）实现JWT Token的签发和解析。Token结构为标准的三段式JWT格式：\verb|Header.Payload.Signature|。

Payload（Claims）包含以下字段：

\begin{table}[htbp]
    \centering
    \caption{JWT Token Payload字段定义}
    \begin{tabular}{|c|c|c|}
        \hline
        \textbf{字段名} & \textbf{类型} & \textbf{说明} \\
        \hline
        userId & Long & 用户ID，主键值 \\
        \hline
        username & String & 登录用户名 \\
        \hline
        role & String & 角色标识（ADMIN/STUDENT/TEACHER） \\
        \hline
        iat & Date & 签发时间（Issued At） \\
        \hline
        exp & Date & 过期时间（Expiration），iat + 24小时 \\
        \hline
    \end{tabular}
\end{table}

\subsubsection{Token签名算法}

Token签名采用HMAC-SHA256对称密钥算法，密钥通过\verb|application.properties|中的\verb|jwt.secret|配置项注入。密钥长度建议不少于256位（32字节），以确保足够的签名强度。当前开发环境的密钥为占位符，生产部署前须替换为随机生成的强密钥。

\begin{lstlisting}[language=Java,caption={JwtUtil.generateToken() 核心逻辑}]
public String generateToken(Long userId, String username, String role) {
    Map<String, Object> claims = new HashMap<>();
    claims.put("userId", userId);
    claims.put("username", username);
    claims.put("role", role);

    return Jwts.builder()
            .claims(claims)
            .issuedAt(new Date())
            .expiration(new Date(System.currentTimeMillis() + expiration))
            .signWith(getKey())   // HMAC-SHA256签名
            .compact();
}
\end{lstlisting}

\subsubsection{Token解析与校验}

Token解析时，JJWT库自动校验签名和过期时间。若签名不匹配、Token已过期或格式错误，JwtUtil.parseToken()抛出相应异常，由JwtFilter捕获并统一处理（认证失败时不阻塞请求，仅不注入用户信息，由后续RoleInterceptor按需拒绝）。

\begin{lstlisting}[language=Java,caption={JwtUtil.parseToken() 核心逻辑}]
public Claims parseToken(String token) {
    return Jwts.parser()
            .verifyWith(getKey())
            .build()
            .parseSignedClaims(token)
            .getPayload();
}
\end{lstlisting}

\subsection{请求认证过滤器（JwtFilter）}

JwtFilter是认证机制的核心组件，继承自Spring Web的OncePerRequestFilter抽象类，确保每次请求仅执行一次过滤逻辑。

\begin{lstlisting}[language=Java,caption={JwtFilter.doFilterInternal() 核心逻辑}]
@Override
protected void doFilterInternal(HttpServletRequest request,
        HttpServletResponse response, FilterChain chain)
        throws ServletException, IOException {

    String path = request.getRequestURI();

    // 放行公开端点
    if (path.startsWith("/api/v1/auth/") ||
            path.contains("/doc.html") ||
            path.contains("/v3/api-docs") ||
            path.contains("/swagger") ||
            path.contains("/webjars")) {
        chain.doFilter(request, response);
        return;
    }

    // 解析Token并注入用户信息
    String header = request.getHeader("Authorization");
    if (header != null && header.startsWith("Bearer ")) {
        String token = header.substring(7);
        try {
            if (!jwtUtil.isTokenExpired(token)) {
                Claims claims = jwtUtil.parseToken(token);
                request.setAttribute("userId",
                    claims.get("userId", Long.class));
                request.setAttribute("username",
                    claims.get("username", String.class));
                request.setAttribute("role",
                    claims.get("role", String.class));
            }
        } catch (Exception ignored) {
            // Token无效时不注入属性，后续鉴权拦截器按需拒绝
        }
    }

    chain.doFilter(request, response);
}
\end{lstlisting}

关键设计决策说明：
\begin{itemize}
    \item \textbf{不阻断未认证请求}：Filter不在缺少Token或Token无效时直接返回401，而是将请求继续传递给后续拦截器。这样设计的好处是：若后续Controller方法未标注@RoleRequired，则允许未登录用户访问（为未来可能的公开接口预留灵活性）；
    \item \textbf{白名单路径}：注册、登录、API文档（Swagger/Knife4j）相关路径被明确排除在认证拦截之外；
    \item \textbf{无状态设计}：Filter不查询数据库，仅依赖JWT Token自包含的用户信息完成身份识别，减轻数据库压力。
\end{itemize}

\subsection{角色权限拦截器（RoleInterceptor）}

RoleInterceptor实现Spring MVC的HandlerInterceptor接口，在Controller方法执行前校验用户角色。

\begin{lstlisting}[language=Java,caption={RoleInterceptor.preHandle() 核心逻辑}]
@Override
public boolean preHandle(HttpServletRequest request,
        HttpServletResponse response, Object handler) {

    if (!(handler instanceof HandlerMethod hm)) {
        return true;
    }

    // 获取方法上的 @RoleRequired 注解
    RoleRequired annotation = hm.getMethodAnnotation(RoleRequired.class);
    if (annotation == null) {
        annotation = hm.getBeanType().getAnnotation(RoleRequired.class);
    }
    if (annotation == null) {
        return true;  // 无注解则放行
    }

    String role = (String) request.getAttribute("role");
    if (role == null) {
        throw new BusinessException(401, "请先登录");
    }

    List<String> allowed = Arrays.asList(annotation.value());
    if (!allowed.contains(role)) {
        throw new BusinessException(403,
            "权限不足，需要角色: " + String.join(", ", allowed));
    }

    return true;
}
\end{lstlisting}

拦截器的工作流程如下：
\begin{enumerate}
    \item 判断当前handler是否为Controller方法（HandlerMethod），非方法类型（如静态资源请求）直接放行；
    \item 优先从方法上获取@RoleRequired注解，若不存在则从Controller类上获取，两级均无则放行（无权限要求）；
    \item 从Request Attribute中取出role值（由JwtFilter预先注入），若为null则用户未登录，抛出BusinessException(401)；
    \item 校验当前用户角色是否在@RoleRequired声明的允许角色列表中，不在则抛出BusinessException(403)；
    \item 校验通过后返回true，请求继续进入Controller方法处理。
\end{enumerate}

\subsection{@RoleRequired 注解}

自定义注解@RoleRequired用于在Controller类或方法级别声明接口所需的角色权限。

\begin{lstlisting}[language=Java,caption={@RoleRequired 注解定义}]
@Target({ElementType.METHOD, ElementType.TYPE})
@Retention(RetentionPolicy.RUNTIME)
public @interface RoleRequired {
    String[] value() default {};
}
\end{lstlisting}

典型使用方式：

\begin{itemize}
    \item \textbf{类级别声明}：在Controller类上标注@RoleRequired，该类下所有方法均需指定角色才能访问。例如管理员Controller标注\verb|@RoleRequired("ADMIN")|；
    \item \textbf{方法级别声明}：在特定方法上标注，覆盖类级别设置。例如在学生Controller中，大部分方法需要STUDENT角色，但某个查询接口允许TEACHER查看，则在方法上标注\verb|@RoleRequired({"STUDENT", "TEACHER"})|。
\end{itemize}

\subsection{拦截器注册配置（WebConfig）}

WebConfig实现WebMvcConfigurer接口，将RoleInterceptor注册到Spring MVC拦截器链中。

\begin{lstlisting}[language=Java,caption={WebConfig 拦截器注册}]
@Configuration
@RequiredArgsConstructor
public class WebConfig implements WebMvcConfigurer {

    private final RoleInterceptor roleInterceptor;

    @Override
    public void addInterceptors(InterceptorRegistry registry) {
        registry.addInterceptor(roleInterceptor)
                .addPathPatterns("/api/**")
                .excludePathPatterns(
                    "/api/v1/auth/**",   // 认证接口公开
                    "/doc.html",          // Knife4j文档页
                    "/v3/api-docs/**",    // OpenAPI JSON
                    "/swagger/**",        // Swagger静态资源
                    "/webjars/**"         // Webjars资源
                );
    }
}
\end{lstlisting}

关键配置说明：
\begin{itemize}
    \item \verb|addPathPatterns("/api/**")|：拦截所有API请求；
    \item \verb|excludePathPatterns(...)|：排除认证接口和API文档相关路径，这些路径在JwtFilter中同样被放行，形成双重保障。
\end{itemize}

\subsection{密码加密方案}

系统采用BCrypt单向哈希算法存储用户密码。BCrypt是专门为密码哈希设计的算法，具有以下优势：

\begin{itemize}
    \item \textbf{内置盐值}：每次加密自动生成随机盐值并嵌入哈希结果中（格式：\verb|$2a$10$<22位盐值><31位哈希值>|），相同密码每次加密结果不同，有效防御彩虹表攻击；
    \item \textbf{计算代价可调}：work factor参数（示例为10）表示$2^{10}$轮哈希迭代，可根据硬件性能调整。当前默认值10在保证安全性的同时，单次加密耗时约100ms，对登录体验影响可忽略；
    \item \textbf{不可逆}：BCrypt是单向哈希，无法从密文反推明文密码。
\end{itemize}

AuthService中直接创建BCryptPasswordEncoder实例：

\begin{lstlisting}[language=Java]
private final BCryptPasswordEncoder passwordEncoder =
    new BCryptPasswordEncoder();
\end{lstlisting}

注册时调用\verb|passwordEncoder.encode(rawPassword)|生成密文存储；登录时调用\verb|passwordEncoder.matches(rawPassword, encodedPassword)|比对密码。

\subsection{Redis Token缓存策略}

Token在签发的同时写入Redis缓存，key设计为\verb|"token:" + userId|，value为Token字符串本身。设置过期时间为24小时，与JWT Token的过期时间保持一致。

Redis Token缓存的作用：
\begin{itemize}
    \item \textbf{登出管理}：前端登出时，后端删除Redis中的Token记录，使得该Token不再被认可（即便JWT本身未过期）；
    \item \textbf{单点登录控制}：可选扩展——同一userId的Token写入Redis时覆盖旧值，天然实现了同一账号后续登录踢掉前一次登录的效果；
    \item \textbf{状态查询}：管理员可通过Redis中是否存在某用户的Token来判断其在线状态。
\end{itemize}

RedisTemplate的序列化配置（RedisConfig）采用String序列化key、GenericJackson2JsonRedisSerializer序列化value，支持对象数据的JSON序列化存储，同时保持Redis客户端中的key可读性。

\subsection{统一返回体设计}

所有API响应均通过Result<T>泛型类统一封装，结构如下：

\begin{lstlisting}[language=Java,caption={Result.java 统一返回体}]
@Data
@NoArgsConstructor
@AllArgsConstructor
public class Result<T> {
    private int code;      // 业务状态码
    private String message; // 提示信息
    private T data;        // 响应数据（泛型）

    public static <T> Result<T> ok(T data) {
        return new Result<>(200, "success", data);
    }

    public static <T> Result<T> error(int code, String message) {
        return new Result<>(code, message, null);
    }
}
\end{lstlisting}

状态码约定：
\begin{itemize}
    \item 200：请求成功；
    \item 400：参数校验失败；
    \item 401：未登录或Token无效；
    \item 403：角色权限不足；
    \item 500：服务端内部错误。
\end{itemize}

\subsection{全局异常处理}

GlobalExceptionHandler通过@RestControllerAdvice注解拦截Controller层抛出的异常，转换为统一返回体格式：

\begin{itemize}
    \item \textbf{BusinessException}：业务异常，提取异常中的code和message构造Result返回；
    \item \textbf{MethodArgumentNotValidException}：JSR-303参数校验异常，汇总所有字段级别的错误信息，以分号分隔返回；
    \item \textbf{Exception}：兜底异常处理，记录完整堆栈日志，向前端返回500错误。
\end{itemize}

该机制确保所有异常（包括认证鉴权模块抛出的BusinessException）都能以统一的JSON格式返回前端，前端只需关注Result中的code和message字段即可完成统一的错误处理逻辑。

\subsection{安全防护措施}

认证权限模块在设计中考虑了以下安全防护点：

\begin{enumerate}
    \item \textbf{密码不落盘}：密码仅以BCrypt密文形式存储在数据库中，日志中不打印密码参数（AuthService中密码字段仅参与比对，不输出日志）；
    \item \textbf{Token时效性}：JWT Token有效期24小时，过期后必须重新登录，降低Token泄露后的安全风险窗口；
    \item \textbf{接口级权限控制}：每个需要权限的API接口通过@RoleRequired明确声明角色要求，防止水平越权（学生调用教师接口）和垂直越权（普通用户调用管理员接口）；
    \item \textbf{注册角色限制}：AuthService.register()方法中硬编码禁止注册ADMIN角色，管理员账号仅由DBA在数据库层面手动创建，防止恶意注册管理员；
    \item \textbf{账号禁用机制}：enabled字段支持管理员禁用违规或异常账号，禁用账号在登录时被明确拒绝（返回403禁用提示，而非模糊的401认证失败）。
\end{enumerate}

\subsection{认证鉴权全链路时序}

用户从登录到访问受保护API的完整认证鉴权链路如下：

\begin{enumerate}
    \item 用户通过\verb|POST /api/v1/auth/login|提交用户名和密码；
    \item AuthService验证用户名密码，通过后生成JWT Token，写入Redis，返回LoginResponse（包含Token）；
    \item 前端将Token存储在localStorage中，并设置Axios请求拦截器——每次HTTP请求自动在header中注入\verb|Authorization: Bearer <token>|；
    \item 用户访问受保护的API（如\verb|GET /api/v1/thesis/list|）：
    \begin{enumerate}
        \item[(a)] JwtFilter从header提取Token，校验签名和过期时间，将userId/username/role注入Request Attribute；
        \item[(b)] RoleInterceptor检查目标Controller方法上的@RoleRequired注解，校验当前用户角色是否在允许列表中；
        \item[(c)] 鉴权通过后，Controller方法通过request.getAttribute("userId")获取当前操作用户ID，确保数据隔离（学生只能操作自己的论文，教师只能查看自己学院的论文等）。
    \end{enumerate}
    \item 用户登出时调用\verb|POST /api/v1/auth/logout|，服务端删除Redis中的Token记录，前端清空localStorage中的Token。
\end{enumerate}

该链路的设计确保了从认证到鉴权的完整闭环，每个环节都有明确的职责边界和异常处理策略，构成了系统安全体系的基石。
